\documentclass[11pt]{article}
\usepackage[a4paper,margin=25mm]{geometry}
\usepackage{amsmath,amssymb,amsthm,hyperref}
\newtheorem{theorem}{Theorem}\newtheorem{lemma}[theorem]{Lemma}
\newcommand{\C}{\mathbb C}\newcommand{\R}{\mathbb R}
\newcommand{\norm}[1]{\left\lVert#1\right\rVert}
\title{The heat depth and the fixed-packet attained limit}
\author{Split-Zero research programme}\date{19 September 2026}
\begin{document}\maketitle
This working continuation of HM1--30 keeps the original fixed divisor
$h$, its full jet space $E_h$, canonical section $R_h$ and physical
Hilbert space $L^2(\R_{>0},dx)$. It determines the dependence of the
proved heat comparison on the actual finite metric. The underlying
source-density calculation was recovered as RC6--10 in
\texttt{EXACT\_CORPUS\_RECEIVERS.tex}; its needed proof is reproduced
here. The accepted HM source is preserved at GitHub commit
\texttt{88f967d3c38750c57264fc5ebbf10bb463717e4e}.

\section{Density in the stated physical norm}
Put $D=-x\partial_x$, $f_j=D^j\Theta\phi_*$ and
$g_0(t)=2\xi(1/2+it)$, with the original $\phi_*$ of HM1.
Mellin Parseval sends $f_j$ to $(1/2+it)^jg_0(t)$ in
$L^2(\R,dt/(2\pi))$. We first establish the exact density needed
for the finite-metric conclusion.

The bounded alternating partial sum
$A(y)=\sum_{n\le y}(-1)^{n-1}$ gives
$\eta(s)=s\int_1^\infty A(y)y^{-s-1}dy$ for $\Re s>0$.
On $s=1/2+it$ this has absolute value at most $2|s|$.
Continuation of $\eta(s)=(1-2^{1-s})\zeta(s)$ from $\Re s>1$
then gives $|\zeta(s)|\le2|s|/(\sqrt2-1)$ on this line.
Stirling's formula in its fixed vertical strip, with the full Gamma
factor in $2\xi$, yields
\begin{equation}\tag{HD1}
 |g_0(t)|\le C(1+|t|)^{11/4}e^{-\pi|t|/4},\qquad
 \int_\R |g_0(t)|^2e^{2a|t|}dt<\infty\quad(0<a<\pi/4).
\end{equation}
The Gamma asymptotic is the classical Stirling formula cited in
\cite{gamma}; the partial-summation bound was just proved.

If $v$ is orthogonal to all the polynomial multiples of $g_0$, put
$b(t)=\overline{v(t)}g_0(t)$. Cauchy--Schwarz and (HD1) give
$b\in L^1(e^{a|t|}dt)$ for every $a<\pi/4$, and all its polynomial
moments vanish. Its bilateral transform
$B(z)=\int b(t)e^{zt}dt/(2\pi)$ is holomorphic on $|\Re z|<\pi/4$:
on every smaller closed strip all derivatives are dominated using a
slightly larger $a$ in (HD1). Its Taylor coefficients at zero vanish,
so the identity theorem gives $B=0$ on this strip. Fourier uniqueness
for $L^1$ functions \cite{Fourier} implies $b=0$ almost everywhere.
The zeros of the nonzero entire function $\xi$ are discrete, so
$g_0\ne0$ almost everywhere and $v=0$. We have proved
\begin{equation}\tag{HD2}
 \overline{\operatorname{span}\{f_j:j\ge0\}}^{L^2(dx)}=L^2(dx).
\end{equation}
This proof uses no assumption about the real parts of the zeros.

\section{The exact finite constant and the required heat scale}
Fix the complete nonempty divisor $h$ throughout this section and
write $d=\deg h$. For $N\ge d$, retain the matrices of HM10:
\[
 \mathbb G_N=\begin{pmatrix}A&C_N\\C_N^*&S_N\end{pmatrix},\quad
 A=R_h^*R_h>0,\quad H_N=A-C_NS_N^{-1}C_N^*>0.
\]
The relation columns are exactly $f_0,\ldots,f_{N-d}$; the canonical
section and $A$ are independent of $N$. In particular every statement
below fixes the divisor rather than replacing it by a growing one.

\begin{theorem}
The full comparison constant in HM10 is exactly
\begin{equation}\tag{HD3}
 c_N^2=\lambda_{\max}(H_N^{-1/2}AH_N^{-1/2})
      =\sup_{u\ne0}\frac{\norm{R_hu}_{L^2(dx)}^2}{u^*H_Nu}.
\end{equation}
Moreover $H_N$ decreases in Loewner order to zero in operator norm,
and $c_N$ is nondecreasing for $N\ge d$ with $c_N\to\infty$.
For the literal heat replacement in HM9, set
$\delta=2^{-(J+1)/2}c_N$. Then
\begin{equation}\tag{HD4}
 \lambda_{\max}\!\left(\mathbb G_N^{-1/2}
                \mathbb G_N^{[J]}\mathbb G_N^{-1/2}\right)
       \ge (\delta-1)_+^2.
\end{equation}
Thus any fixed finite upper comparison constant $B>0$ for these full
prequotients requires
\begin{equation}\tag{HD5}
 J+1\ \ge\ 2\log_2 c_N-2\log_2(1+\sqrt B).
\end{equation}
The sufficient depth of HM24 and this necessary depth therefore have
the same leading dependence $2\log_2 c_N$, with their stated fixed
tolerance constants retained. No rate for $c_N$ in $N$ is asserted.
\end{theorem}
\begin{proof}
Since $C_NS_N^{-1}C_N^*=A-H_N$, the definition in HM10 becomes
\[
 c_N^2=1+\lambda_{\max}
       (H_N^{-1/2}(A-H_N)H_N^{-1/2})
       =\lambda_{\max}(H_N^{-1/2}AH_N^{-1/2}).
\]
Substituting $v=H_N^{1/2}u$ in its Rayleigh quotient proves the last
formula in (HD3), without changing the original coefficient maps.

The space of relation columns increases with $N$, so its attained
distance form $H_N$ decreases. The density (HD2) makes the distance
from every column of $R_h$ to this relation space tend to zero.
Hence every diagonal entry of $H_N$ tends to zero. Positivity bounds
the absolute value of every off-diagonal entry by the geometric mean
of the corresponding diagonal entries. The fixed dimension is finite,
so $\norm{H_N}\to0$. At each finite $N$ it is positive definite:
if $R_hu$ were in the finite relation span, the full original Mellin
jet map would give $u=0$, since it kills all the $f_j$.
Formula (HD3) proves monotonicity of $c_N$. It also gives
$c_N^2\ge\lambda_{\min}(A)/\norm{H_N}\to\infty$.

Let $\mathcal P_N=[0,F_{N-d}]$ and $s=2^{J+1}$, as in HM7--9.
The operator $U_N=\mathcal R_N\mathbb G_N^{-1/2}$ is an isometric
embedding, so $\norm{U_N}=1$. By the exact computation in HM13 and
$T_s^*T_s=s^{-1}I$,
\[
 \norm{T_s\mathcal P_N\mathbb G_N^{-1/2}}=s^{-1/2}c_N=\delta.
\]
The reverse triangle inequality implies
$\norm{\mathcal R_N^{[J]}\mathbb G_N^{-1/2}}\ge(\delta-1)_+$.
Squaring gives (HD4). If the left side of (HD4) is at most $B$,
then $\delta\le1+\sqrt B$, which is exactly (HD5). At fixed $J$
the lower bound in (HD4) tends to infinity. Every map here is the
full original prequotient map with the same section and coefficient
fibre as HM6.
\end{proof}

\section{The receiving quotient is still taken through its original map}
The lower bound (HD4) concerns the full form before minimization.
Its exact receiving map is the already proved
$M\mapsto(\Lambda M^{-1}\Lambda^*)^{-1}$. A lower bound on the
largest eigenvalue of a full form need not survive that map. For an
explicit verification, let $M_t=\operatorname{diag}(t,1)$ on $\C^2$,
$M_0=I$, and $\Lambda(z_1,z_2)=z_2$. Then
$\lambda_{\max}(M_t)=t$ for $t\ge1$, but
$\Lambda M_t^{-1}\Lambda^*=1$ and the attained quotient form is
identically $1$. The fixed-divisor limit above therefore proves the
stated need to increase heat depth for the full-prequotient comparison;
it does not assign a relative-rate failure to the re-minimized packet
metrics. Those continue to obey the explicit two-sided bounds of
HM19 whenever the depth in HM24 is used.

The exact representatives
$R_hu-F_{N-d}S_N^{-1}C_N^*u$ keep the original full jets and theta
cohomology class while their squared norms are $u^*H_Nu\to0$.
Thus the inclusion of the original rapid-decay source space into
$L^2(dx)$ induces the zero map to the Hausdorff quotient by the
closed theta image. Its finite jet sections still split the original
algebraic cohomology. These explicit maps preserve the distinction
between the chosen source topology and its Hilbert completion, and
leave the original growing-packet arithmetic estimates intact.

\begin{thebibliography}{9}
\bibitem{gamma} Richard A.~Askey and Ranjan Roy,
NIST Digital Library of Mathematical Functions,
Chapter 5, \S5.11, Stirling's formula,
\url{https://dlmf.nist.gov/5.11.E3}.
\bibitem{Fourier} Elias M.~Stein and Rami Shakarchi,
\emph{Fourier Analysis: An Introduction}, Princeton University Press,
2003, Chapter 5: Fourier inversion and Plancherel on the real line.
\end{thebibliography}
\end{document}
